%! Equivalent Representations of Polynomial Expressions — Unit 2, Lesson 2.6 — Lesson Plan
\documentclass[10pt]{article}
\usepackage{precalculus-article}
\usepackage{precalculus-boxes}
\newcommand{\TallMath}[1]{$\displaystyle #1\rule[-1.4em]{0pt}{3.2em}$}
\newcommand{\CourseName}{Precalculus}
\newcommand{\MeetingLength}{60 minutes}
\newcommand{\UnitNumberName}{Unit 2: Polynomial Functions \quad}
\newcommand{\LessonNumberName}{Lesson 2.6: Equivalent Representations of Polynomial Expressions}

\begin{document}
\begin{center}
    {\Large\bfseries \CourseName \\
    {\small\bfseries \UnitNumberName \LessonNumberName}}
\end{center}

\begin{tcolorbox}[colback=lilac, colframe=plum, boxrule=0.9pt, arc=2mm,
    left=2mm, right=2mm, top=1.5mm, bottom=1.5mm]
    \textbf{Primary Objective:} Students will decide which analytic form of a
    polynomial --- factored or standard --- answers a given question, convert
    between the two by multiplying out and by using the binomial theorem with
    Pascal's Triangle, carry out polynomial long division to write
    $f(x) = g(x)q(x) + r(x)$ with $\deg r < \deg g$, and use both forms of one
    polynomial to answer questions in an applied context.
    \hfill\textit{\small (CED Topic 1.11, polynomial half --- Skills 1.B, 3.B)}
\end{tcolorbox}

\renewcommand{\arraystretch}{1.6}
\begin{skillbox}[Priority Ideas \& Skills]{goldbox}
\begin{tabularx}{\linewidth}{@{}X X@{}}
\textbf{AP Skills} &
\textbf{Key Understandings} \\
\midrule
\textbf{1.B} --- Express functions, equations, or expressions in analytically
equivalent forms useful in a given context. &
\begin{itemize}[leftmargin=1.2em,itemsep=2pt,topsep=0pt]
  \item The \textbf{factored} form readily provides the real zeros, so it
        reveals the $x$-intercepts and the intervals that bound domain and
        range questions in context. (1.11.A.1)
  \item The \textbf{standard} form displays the leading term, so it reveals
        end behavior --- and its constant term is the $y$-intercept.
        (1.11.A.2)
  \item The binomial theorem uses one row of Pascal's Triangle to expand
        $(a+b)^n$, including $p(x) = (x+c)^n$. (1.11.C.1)
\end{itemize} \\
\textbf{3.B} --- Apply numerical results in a given mathematical or applied
context. &
\begin{itemize}[leftmargin=1.2em,itemsep=2pt,topsep=0pt]
  \item Polynomial long division mirrors numerical long division: dividing $f$
        by $g$ gives $f(x) = g(x)q(x) + r(x)$ with $\deg r < \deg g$.
        (1.11.B.1)
  \item Remainder $0$ is exactly the statement that $g$ is a \textit{factor}
        of $f$ --- the factor/zero equivalence from Lesson 2.3, now with a
        procedure attached to it.
  \item Information pulled from different forms of the \textit{same}
        polynomial answers different questions about one situation.
        (1.11.A.3)
\end{itemize} \\
\end{tabularx}
\end{skillbox}

\begin{skillbox}[{Vocabulary, Concepts, \& Theorems}]{greenbox}
\renewcommand{\arraystretch}{1.4}
\begin{tabularx}{\linewidth}{@{} >{\leavevmode\bfseries\color{plum}}p{0.27\linewidth} X @{}}
Standard form     & A polynomial written as a sum of terms in descending degree, $a_nx^n + \cdots + a_1x + a_0$. The leading term is visible, and $a_0$ is the $y$-intercept. \\
Factored form     & The same polynomial written as a product of factors. Each linear factor $(x - a)$ hands over the real zero $a$. \\
Pascal's Triangle & The number triangle whose row $n$ lists the coefficients of $(a+b)^n$; each entry is the sum of the two directly above it. \\
Binomial theorem  & The rule that expands $(a+b)^n$ using row $n$ of Pascal's Triangle, with the power of $a$ falling and the power of $b$ rising across the terms. \\
Polynomial long division & The algorithm that divides one polynomial by another, one leading term at a time, exactly as numerical long division does. \\
Quotient and remainder & In $f(x) = g(x)q(x) + r(x)$, $q$ is the quotient and $r$ is the remainder. \\
Division algorithm & For polynomials $f$ and $g \ne 0$ there are $q$ and $r$ with $f(x) = g(x)q(x) + r(x)$ and $\deg r < \deg g$. \\
\end{tabularx}
\end{skillbox}

\begin{skillbox}[Activate Prior Knowledge \& Spiral Review]{lilac}
    Please see attached warm-up (spiral review).
    Skills reviewed: reading degree, leading coefficient, and end behavior in
    limit notation off a polynomial in standard form (Lesson 2.5); listing the
    real zeros with their multiplicities off a polynomial in factored form
    (Lesson 2.3); and one Algebra 2 item multiplying two binomials. Items 1 and
    2 are deliberately the two forms this lesson compares --- each answers a
    question the other cannot, which is the whole lesson in ninety seconds.
\end{skillbox}

\begin{skillbox}[Hook]{lilac}
    \textbf{One function, two costumes.} Put both lines on the board at once
    and say nothing about which is which:
    \begin{enumerate}[itemsep=2pt]
        \item $p(x) = (x+1)(x-2)(x-3)$ \qquad and \qquad
              $p(x) = x^3 - 4x^2 + x + 6$.
        \item Ask groups to verify with one multiplication that these really
              are the same function --- then ask two questions:
              \textit{Where does the graph cross the $x$-axis?} and
              \textit{What does the graph do as $x \to -\infty$?}
    \end{enumerate}
    \textit{Discussion prompt:} nobody reaches for the same line twice. The
    intercepts fall out of the product; the end behavior falls out of the
    leading term $x^3$. Neither form is ``simplified'' and neither is
    ``better'' --- today we learn to \textit{pick} the form the question wants,
    and to move between them on demand.
\end{skillbox}

\begin{skillbox}[Lesson]{lilac}
\begin{multicols}{2}

\textbf{Part 1: Which Form Answers Which Question?} (12 min --- the spine)

\begin{itemize}
  \item Build the two-column table with the class. Factored form: real zeros
        (2.3), hence $x$-intercepts, hence the intervals that matter in
        context. Standard form: leading term, hence end behavior (2.5), and
        the constant term, hence the $y$-intercept.
  \item Degree is the one thing \textit{both} forms give: add the exponents on
        the factors, or read the highest power.
  \item Run the hook polynomial through every row of the table so students see
        one function answering four questions from two forms.
  \item Say the AP language out loud: this is choosing an equivalent form
        \textit{useful in a given context} --- Skill 1.B.
\end{itemize}

\textbf{Part 2: Factored $\rightarrow$ Standard, and the Binomial Theorem}
(12 min)

\begin{itemize}
  \item Expanding is just distributing, two factors at a time. Model
        $(x+1)(x-2)(x-3)$: pair the first two, then multiply by the third.
  \item Then the special case worth a tool: $(x+c)^n$. Multiplying four copies
        of $(x+2)$ by hand is slow and error-prone.
  \item Show Pascal's Triangle rows 0--5 already printed. Row $n$ gives the
        coefficients of $(a+b)^n$, and each entry is the sum of the two
        directly above it.
  \item Expand $(x+2)^4$ with row 4 ($1, 4, 6, 4, 1$): powers of $x$ fall
        $4,3,2,1,0$ while powers of $2$ rise $0,1,2,3,4$. Result
        $x^4 + 8x^3 + 24x^2 + 32x + 16$. Keep this brisk --- it is a tool, not
        a derivation.
\end{itemize}

\columnbreak

\textbf{Part 3: Long Division and the Division Algorithm} (19 min --- the
core)

\begin{itemize}
  \item Open with the numerical analogue: $17 \div 5$ is $3$ with remainder
        $2$, so $17 = 5 \cdot 3 + 2$, and the remainder is smaller than the
        divisor.
  \item State the polynomial version: $f(x) = g(x)q(x) + r(x)$ with
        $\deg r < \deg g$. ``Smaller'' now means \textit{lower degree}.
  \item Model the tableau for $2x^3 - 3x^2 - 5x + 6$ divided by $x-2$: divide
        leading by leading, multiply, subtract, bring down. Remainder $0$,
        quotient $2x^2 + x - 3$.
  \item Now cash it in against Lesson 2.3: remainder $0$ $\Longleftrightarrow$
        $(x-2)$ is a \textbf{factor} $\Longleftrightarrow$ $2$ is a zero. So
        long division is how you \textit{test} a proposed factor and how you
        \textit{reduce} a polynomial once one zero is known.
  \item Second model, with a nonzero remainder: $x^3 - 4x^2 + 2x + 5$ divided
        by $x-3$ gives $q(x) = x^2 - x - 1$ and $r = 2$. Write the
        division-algorithm sentence and check the degree condition aloud.
\end{itemize}

\textbf{Part 4: Both Forms, One Context} (10 min)

\begin{itemize}
  \item Open box cut from a $12 \times 16$ sheet:
        $V(x) = x(12-2x)(16-2x) = 4x^3 - 56x^2 + 192x$.
  \item Factored form answers ``for what $x$ is the volume $0$?'' and so pins
        the physical domain $0 < x < 6$. Standard form answers ``what is the
        leading term, and what does the model do for large $x$?''
  \item Land Skill 3.B: the numbers mean something here. $x = 8$ is an
        algebraic zero the box cannot have.
\end{itemize}
\end{multicols}
\end{skillbox}

\begin{skillbox}[Explicit Instruction --- The Long-Division Routine]{lilac}
Model the same four beats every single time, and name them aloud while you
write:
\begin{enumerate}[itemsep=2pt]
  \item \textbf{Divide:} leading term of what is left $\div$ leading term of
        the divisor. That is the next term of the quotient.
  \item \textbf{Multiply:} that quotient term times the whole divisor.
  \item \textbf{Subtract:} the leading terms must cancel. If they do not, the
        divide step was wrong.
  \item \textbf{Bring down:} the next term, and repeat until what is left has
        lower degree than the divisor. What is left is $r(x)$.
\end{enumerate}
Insist on placeholder zeros for missing degrees --- $x^3 + 3x^2 - 2$ is
written $x^3 + 3x^2 + 0x - 2$. A missing column is where most wrong answers
start.
\end{skillbox}

\begin{skillbox}[Active Monitoring]{redbox}
While students work on guided notes and practice, circulate and check:
\begin{itemize}
  \item \textbf{Common error:} reaching for the factored form to get end
        behavior. Cold-call: ``Which term wins when $x$ is huge? Can you see
        that term in this form?''
  \item \textbf{Common error:} calling standard form ``simplified'' and
        treating the conversion as one-directional. Both forms are equally
        finished; ask ``which one does \textit{this} question want?''
  \item \textbf{Common error:} in the subtraction step, distributing the minus
        sign to only the first term. Require the subtracted line in
        parentheses before subtracting.
  \item \textbf{Common error:} dropping a placeholder $0x$ column, which
        shifts every later term by one degree.
  \item \textbf{Common error:} applying the Pascal row but forgetting to raise
        $c$ to a power --- $(x+2)^4$ coming out as
        $x^4 + 8x^3 + 12x^2 + 8x + 2$.
  \item \textbf{Language:} ``quotient'' and ``remainder,'' not ``the answer''
        and ``the leftover.''
\end{itemize}
\end{skillbox}

\begin{skillbox}[Group Work \& Differentiation]{redbox}
Students work on the \textbf{Group Activity: Same Polynomial, Two Jobs} in
leveled tiers. Every tier works from printed expressions --- no tier draws
anything.

\begin{itemize}
  \item \textbf{Tier R (Remediate):} Label four questions ``factored'' or
        ``standard,'' complete row 3 of Pascal's Triangle and use it on
        $(x+1)^3$, expand one pair of binomials, and complete the
        division-algorithm sentence.
  \item \textbf{Tier A (Approaching Proficiency):} Take
        $p(x) = (x+2)(x-1)(x-3)$ through both forms and pull four different
        facts out of them, expand $(x+3)^4$ with row 4, and divide
        $x^3 + 2x^2 - 5x - 6$ by $x+1$.
  \item \textbf{Tier E (Extension):} Test the proposed factor $(x-2)$ against
        $x^3 - 6x^2 + 11x - 6$ by division and factor it completely; divide by
        a \textit{quadratic} divisor ($x^3 + 3x^2 - 2$ by $x^2 - 1$) and check
        the degree condition on the remainder; and justify in writing why
        ``remainder $0$'' and ``the divisor is a factor'' are the same
        statement (AP Skill 1.B).
\end{itemize}
\end{skillbox}

\begin{skillbox}[Individual Work \& Assessment]{redbox}
\textbf{Exit Ticket} (last 5--7 minutes, individual, collected):
\begin{enumerate}
  \item One function in two forms, $q(x) = (x+1)(x-5) = x^2 - 4x - 5$: name
        which form answers each of three questions, then answer them.
  \item Expand $(x+2)^3$ using row 3 of Pascal's Triangle (worked space
        provided).
  \item Divide $x^2 + 5x + 1$ by $x + 2$, then say whether $x+2$ is a factor.
\end{enumerate}
\textbf{Assessment note:} Item 1 is the diagnostic --- a student who writes
``factored'' for the end-behavior question has not separated the two forms,
and that separation is the whole lesson. Item 3's remainder of $-5$ checks
that they read ``not a factor'' off the remainder rather than off the look of
the quotient.
\end{skillbox}

\begin{skillbox}[Reinforcement \& Extension]{goldbox}
\textbf{Homework:} 5 problems covering the form/information decision on a
cubic given in both forms, two expansions into standard form, a Pascal's
Triangle row plus two binomial expansions, one long division of
$x^3 - 2x^2 - 5x + 6$ by $x - 1$ with the division-algorithm sentence and a
complete factorization, and one AP-style multiple-choice item on a remainder.

\textbf{Extension (optional):} Substitute $x = a$ into
$f(x) = (x-a)q(x) + r$ --- every term carrying the factor $(x-a)$ dies, so
$r = f(a)$. That is the remainder theorem, and it turns Lesson 2.3's factor
theorem into a one-evaluation test.

\textbf{Preview:} Lesson 2.7, \textit{Polynomial Equations, Inequalities, and
Modeling}, puts these forms to work: solving $p(x) = 0$ starts by getting to
factored form, and reading where $p(x) > 0$ starts from the zeros and the end
behavior --- one fact from each form.
\end{skillbox}

\begin{teachernote}[Warm-Up]
Five minutes, silent start, no calculators. Items 1 and 2 are the two halves
of today's argument, so review them side by side rather than in order: ask
whether item 1's form could have answered item 2's question, and the reverse.
It cannot, and that exchange is worth more than either answer. Item 3 is the
Algebra 2 hinge --- a class that cannot multiply two binomials cleanly will
not survive Part 2, so if more than a few students stumble, budget five extra
minutes there and compress Part 4 to the domain question only.
\end{teachernote}

\begin{teachernote}[Guided Notes]
Parts 1 and 3 carry the lesson; Part 2 is real content but it is a tool, so
resist the urge to derive Pascal's Triangle. State the addition rule, point at
one entry, expand $(x+2)^4$, move on. In Part 1, fill the form/information
table \textit{with} the class rather than handing it over --- the table is the
lesson's one-page summary, and students who copy it passively still reach for
the wrong form on the exit ticket. Part 3 needs the full nineteen minutes:
write the first tableau slowly, narrating divide--multiply--subtract--bring
down, and do the subtraction in parentheses on the board every time. The
payoff line is the one connecting back to 2.3 --- say it explicitly, ``the
remainder is zero exactly when the divisor is a factor,'' and write the double
arrow. If the clock runs out, cut Part 4's second question, never the second
division example: a class that has only ever seen a remainder of $0$ has not
seen the division algorithm.
\end{teachernote}

\begin{teachernote}[Group Activity]
Tier R's first item is deliberately not a computation --- it is the sorting
decision, repeated four times, and a group that gets it wrong needs Part 1
retaught before any division practice will stick. Tier A's item 3 is the
standard-issue long division; check the subtraction lines, not the final
answer, because a right answer reached through a wrong subtraction will not
repeat. Tier E's quadratic divisor separates students who learned the routine
from students who memorized ``divide by $x$ minus something,'' and its
remainder $x + 1$ is the first time $\deg r < \deg g$ has any content, since a
linear remainder is not a constant. Item 3 is the writing task; accept any
version of ``if the remainder is zero then $f(x) = g(x)q(x)$, and that is what
factor means.''
\end{teachernote}

\begin{teachernote}[Exit Ticket]
Collected, completion credit only. Sort by item 1's form labels first: a
student who answers all three questions correctly but labels the forms wrong
got there by computing rather than by choosing, and that is a Part 1 reteach
even though the answers are right. On item 3, a remainder of $-5$ paired with
``yes, it is a factor'' is a definition miss and worth naming to the whole
class tomorrow. The ticket should take five minutes; if item 2 alone takes
that long, the binomial theorem needs one more example before 2.7.
\end{teachernote}

\begin{teachernote}[Homework]
Problem 1 is the one to grade first --- it isolates the form/information
decision from every procedure in the lesson. Problem 4 uses
$x^3 - 2x^2 - 5x + 6$, the same polynomial from Lesson 2.3's warm-up and
guided notes, so students should recognize that $x = 1$ was already confirmed
as a zero; if nobody notices, open tomorrow with the two pages side by side.
Problem 5 is the AP-style item --- distractor D (``$0$'') is chosen by
students who assume every division comes out even, which is exactly what the
second guided-notes example exists to prevent. The extension is the remainder
theorem; if a group asked ``do we have to divide every time?'' during the
lesson, assign it explicitly and open 2.7 with a student presenting it.
\end{teachernote}
\end{document}
